problem:  
Let \( (M^n,g) \) be a compact, simply connected Riemannian manifold of dimension \( n \ge 5 \) with sectional curvature \( 0 \le K \le 1 \), and suppose the set of points where \( K>0 \) is open and dense (so \( M \) has almost positive curvature). Assume that a compact Lie group \( G \) acts effectively and isometrically on \( M \) with cohomogeneity one (so the orbit space \( M/G \) is one–dimensional).  

**Question:**  
For fixed \( n \) and \( G \), are there only finitely many diffeomorphism types of such manifolds \( M \)?  

Why is it a "good" problem:  
This problem directly targets the borderline between positive and almost positive curvature under symmetry constraints, an area where very few classification results exist. Cohomogeneity–one actions admit a precise description by a *group diagram* \( H \subset K_\pm \subset G \), which encodes the manifold’s topology. For manifolds with strictly positive curvature, curvature restrictions severely limit these diagrams and suggest finiteness up to diffeomorphism; yet it is unknown whether the same rigidity persists when curvature is merely almost positive.  

Resolving this question would therefore test the robustness of the known rigidity principles under minimal relaxation of curvature positivity. A positive answer would extend finiteness phenomena known in nonnegative and positive curvature (e.g., in Grove–Ziller type constructions) into the more flexible almost–positive regime, while a counterexample would reveal previously unseen topological freedom. The problem beautifully links Riemannian geometry (curvature structure), transformation groups (cohomogeneity–one classification), and differential topology (diffeomorphism type). It is conceptually simple but demands new geometric insights into how curvature, symmetry, and topology constrain one another—capturing the essence of a "narrow entrance, profound depth" research problem.